\documentclass[11pt]{article}
\usepackage{amsmath,amssymb,amsthm,hyperref}
\newtheorem{theorem}{Theorem}
\newtheorem{lemma}[theorem]{Lemma}
\newcommand{\del}{\delta}
\title{Deterministic six-window discrepancy for short permutation families}
\author{Research note; independently reviewed by two other agents}
\date{30 September 2026}
\begin{document}
\maketitle

For a probability distribution $\mu$ supported on $q$ permutations of
$[n]$, define, on distinct labels,
\[
 \del(a,b,c)=\frac12-\frac32
       \Pr_{\rho\sim\mu}(\rho(b)<\rho(a),\rho(b)<\rho(c)).
\]
Thus $|\del|\le1$, $\del(a,b,c)=\del(c,b,a)$, and the sum over the
three choices of center is zero. For a target permutation $\sigma$, put
\[
 F(\sigma)=\sum_{j=1}^{n-2}
       \del(\sigma_j,\sigma_{j+1},\sigma_{j+2}).
\]
This is a constant-position-weight version of the triple contribution
in the palindrome construction. The conclusions below are about this
linear statistic, not the full nonlinear fair-order probabilities.

\begin{theorem}\label{main}
There are absolute constants $C,c>0$ such that, if $n\ge Cq$, some six
distinct labels have an interior middle-pair swap which changes their
four-triple path score by at least $c/\sqrt q$.
If $n\ge2Cq$, then
\[
 \max_\sigma F(\sigma)-\min_\sigma F(\sigma)
       \ge c\frac n{\sqrt q}.
\]
The second conclusion can be witnessed on a family of independent swaps
in disjoint six-label windows.
\end{theorem}

For any six labels define the swap difference
\begin{align*}
 D(a,b,c,d,e,f)={}&\del(a,b,c)+\del(b,c,d)+\del(c,d,e)+\del(d,e,f)\\
 &-\del(a,b,d)-\del(b,d,c)-\del(d,c,e)-\del(c,e,f).
\end{align*}

\begin{lemma}[Small swaps force a pair coboundary]\label{coboundary}
Let $\del$ be any function on distinct triples with the symmetry,
center-sum identity, and bound $|\del|\le1$ stated above. Suppose
$|D(a,b,c,d,e,f)|\le\eta$ for every six distinct labels. Then there
is an antisymmetric array $a_{ij}=-a_{ji}$, with $|a_{ij}|\le3/2$,
such that, for all distinct $b,c,d$,
\[
 |\del(b,c,d)-a_{cb}-a_{cd}|\le\eta+\frac{128}{n}.
\]
\end{lemma}
\begin{proof}
Extend $\del$ to be zero when labels repeat, so all its identities
remain valid. All averages in this proof are uniform over all $n$ labels,
with replacement. Define
\[
 U_b(c)=\mathbb E_a\del(a,b,c),\qquad
 r_b=\mathbb E_c U_b(c),\qquad s_c=\mathbb E_b U_b(c).
\]
Averaging the center-sum identity gives $s_c=-r_c/2$ and
$\mathbb E_c r_c=0$. Set
\[
 K_b(c,d)=\del(b,c,d)-\del(b,d,c)+U_b(c)-U_b(d),
\]
\[
 V(c,d)=\mathbb E_b K_b(c,d)
       =U_c(d)-U_d(c)-r_c/2+r_d/2.
\]
For fixed distinct $b,c,d,e$, averaging $D(a,b,c,d,e,f)$ over $a,f$
gives exactly $K_b(c,d)-K_e(c,d)$. The probability that the six labels
are not distinct is at most $9/n$, and $|D|\le8$. Consequently
\[
 |K_b(c,d)-K_e(c,d)|\le\eta+72/n.
\]
Averaging over $e$, including the at most three exceptional values, and
using $|K_b-K_e|\le8$, yields, for distinct $b,c,d$,
\[
 |K_b(c,d)-V(c,d)|\le\zeta:=\eta+96/n.\tag{2}
\]

For distinct $b,c$, average (2) over $d$. The exceptional two values
contribute at most $14/n$, since $|K_b-V|\le7$. Directly averaging
the definitions gives
\[
 \left|2U_b(c)+2U_c(b)-r_b-r_c\right|
       \le\zeta+14/n=: \zeta_2.\tag{3}
\]
Indeed, $\mathbb E_d[K_b(c,d)-V(c,d)]$ equals the expression inside
the absolute value in (3).
Put
\[
 S_{bc}=U_b(c)+U_c(b)-(r_b+r_c)/2,
 \qquad
 a_{bc}=\frac{U_b(c)-U_c(b)}2+\frac{r_c-r_b}4.
\]
Then $|S_{bc}|\le\zeta_2/2$, $a$ is antisymmetric, and $|a_{bc}|\le3/2$.
Use the center-sum identity to write $3\del(b,c,d)$ as the sum of its
differences with the other two centers. Apply (2) to each difference.
After substituting the definitions, this gives
\[
 3\bigl(\del(b,c,d)-a_{cb}-a_{cd}\bigr)
   =-\tfrac12S_{cb}-\tfrac12S_{cd}+S_{bd}+E,
 \qquad |E|\le2\zeta.
\]
The absolute value is at most $\zeta_2+2\zeta$, proving the stated
bound (in fact $\eta+101/n$ is enough).
\end{proof}

\begin{proof}[Proof of Theorem~\ref{main}]
Let $\eta$ be the largest absolute six-label swap difference, and take
$a$ from Lemma~\ref{coboundary}. Write $m=n-1$,
$\tau=\eta+128/n$, and
$p_{ij}=\Pr(\rho(i)<\rho(j))$. Define
\[
 d_{ij}=p_{ij}-\frac13+\frac43a_{ij}.
\]
Antisymmetry and $p_{ij}+p_{ji}=1$ imply
$d_{ij}+d_{ji}=1/3$. Hence some $i$ satisfies
$\sum_{j\ne i}d_{ij}\ge m/6$.

As before, form the Gram matrix $G$ of the $m$ indicator vectors
$\mathbf1_{\rho(i)<\rho(j)}$ in the weighted $q$-dimensional space.
It has rank at most $q$, diagonal $p_{ij}$, and off-diagonal
\[
 G_{jk}=\frac13-\frac23\del(j,i,k)
       =\frac13-\frac23(a_{ij}+a_{ik})+E_{jk},
 \qquad |E_{jk}|\le\frac23\tau.
\]
Subtract the displayed rank-at-most-two additive matrix, including its
diagonal, to obtain $B=\operatorname{diag}(d_{ij})+E$, where $E$ has
zero diagonal. Then $B$ is symmetric with rank at most $q+2$,
$\operatorname{tr}B\ge m/6$, and $|d_{ij}|\le8/3$. Thus
\[
 \frac m6\le\sqrt{q+2}\,\|B\|_F
 \le\sqrt{q+2}\left(\frac83\sqrt m+\frac23\tau m\right).
\]
It follows that
\[
 \eta\ge\frac1{4\sqrt{q+2}}-\frac4{\sqrt{n-1}}-\frac{128}{n}.
\]
For example, $n\ge2^{16}q+1$ makes the right side at least
$1/(64\sqrt q)$.

The argument is hereditary under restricting to unused labels. Greedily
choose disjoint six-label windows while at least $2^{16}q+1$ labels
remain. In each, permit only the swap of its middle two labels; the
outer two labels on each side stay fixed. Every affected triple lies
entirely in that window. On the resulting binary family of target
permutations, $F$ is therefore a constant plus a sum of independent
binary choices with coefficients $D_j$. Its range is exactly
$\sum_j|D_j|$. For $n$ at least twice a sufficiently large constant
times $q$, there are at least $c n$ windows, proving the result.
\end{proof}

\paragraph{Further refinement and remaining obstacle.}
The companion notes \emph{Weighted collision bounds for triple minima
and six-window swaps} and \emph{A deterministic lower bound for the
binomial triple projection} strengthen this result to weighted effective
support and prove a lower bound of order $n^{13/4}/M^{5/2}$ for the
actual binomial triple contribution whenever $n\le M\le c n^{3/2}$.
These extensions have also been independently reviewed. Control of the
full deterministic nonlinear remainder remains open.
\end{document}
